%=========================================================
% 3. Related Work / Literature Review
% =========================================================
\section*{三、文獻回顧與探討}
\addcontentsline{toc}{section}{三、文獻回顧與探討}

本章節將回顧三個相關研究領域的重要進展：(i) LLM 教育導師之能力與評測、(ii) 長程對話的記憶與跨會議一致性建模、(iii) 人設（persona）表徵與推理時控制（inference-time steering/representation engineering），並說明本研究如何建構於既有成果之上，進一步拓展其理論與方法。

\subsection*{(一) LLM 教育導師：任務形式、能力評測與「長期指導一致性」缺口}
教育導師類系統近年逐漸由「單輪答題」走向「多輪教學互動」與「情境式指導」。BEA 2023 shared task 已將 AI 教師回應生成視為可進行系統化評測的教育對話任務，凸顯教學回應需要兼顧\emph{教學策略、對話連貫與錯誤處理}等要素 \citep{tack-etal-2023-bea}。然而，多數評測仍偏重於短期互動品質，較少直接處理「跨多次會議」的學術指導一致性問題。

為補足評測面向上的缺失，近期開始出現以「導師能力」為核心的系統化評估框架與基準集。例如 Maurya et al.~\citep{maurya-etal-2025-unifying} 提出針對 LLM 導師教學能力的評測架構（taxonomy），將教學能力拆解為多維度（如回饋品質、引導策略、錯誤診斷、適應性等），為本研究後續「指導品質」指標設計提供參考。進一步地，針對 LLM 導師的 benchmark 亦逐漸成形（例如以教學任務、題型或互動策略為主軸），用以量化不同模型與系統設計在教育場景的行為差異 \citep{a2510_02663v1_tutorbench,a2601_21375v1_teachbench}。

另外，除了「回答品質」的評測，亦有研究直接聚焦在「導師—學生」互動的對話策略如何啟動研究討論。例如，Faff~\citep{faff2024prfai_ssrn4999390} 以研究導師主導的提問與回饋策略設計（並結合 PRF-AI 混合工具）探討如何引導學生逐步收斂研究題目與方法路徑。此一脈絡對本研究的啟示在於：虛擬教授不僅需能回答問題，更需在跨會議中持續以「導師式對話策略」推進研究決策與任務分解，才能對應專題指導的真實需求。

此外，人機互動（HCI）研究亦指出學習者與 AI 導師互動時，對「角色穩定性」與「信任」具有高度敏感性：同一系統若在不同回合呈現不一致的指導風格或目標變動，會直接削弱學習者對系統的信賴感與持續使用意願 \citep{tutorup_chi2025,mentigo_chi2025}。綜合而言，現有文獻已提供\emph{教學能力拆解}與\emph{導師基準評測}，但對「長期指導一致性」尚缺少可操作的系統設計與可量測的控制機制；本研究即以此缺口為核心切入點。

\subsection*{(二) 長期對話記憶：從記憶模組到協調策略}
要維持長期一致性，第一個瓶頸在於「記憶」：LLM 若僅依賴上下文視窗或片段摘要，容易在跨會議時遺失研究決策、TODO 與依賴關係。對此，近年多篇工作提出將外部記憶納入 LLM 系統的架構，強調「\emph{短期／長期記憶分工}」與「\emph{檢索＋更新}」迴路。MemoryBank 以長期記憶增強個人化對話，強調對話歷程的持續累積與檢索，以改善後續互動的一致性 \citep{a2305_10250v3_memorybank}。MemGPT 則將 LLM 視為「作業系統式」代理，透過記憶層級與管理策略，在有限上下文下維持長期任務表現 \citep{a2310_08560v2_memgpt}。
同樣地，在高風險領域的個人化助理場景中，亦有研究提出以\emph{短期記憶}（近期對話狀態）與\emph{長期記憶}（穩定使用者偏好／背景）協同運作的設計，並討論兩者在更新頻率、衝突解決與安全性上的權衡 \citep{zhang-etal-2024-llm-based}。近期亦有針對記憶系統的評測與分析工作，指出記憶寫入／淘汰策略、檢索噪聲與錯誤累積，會顯著影響長程對話品質與可靠性 \citep{a2402_17753v1_locomomemoryeval,a2505_16067v2_memorymanagementimpacts}，並提出更系統化的「記憶層管理」觀點與工具化實作 \citep{a2504_19413v1_mem0,a2507_07957v1_mirix}。

然而，上述記憶系統多以「\emph{文本片段}」作為儲存單位，對「研究專題」這類高度依賴\emph{方法假設、實驗依賴與決策鏈}的任務而言，僅靠片段記憶仍可能出現「檢索到資訊卻無法維持推理一致」的情形。本研究因此主張：記憶不只要能保存與取回，還必須可結構化地表徵研究依賴關係，才能在跨會議推理中提供可檢核的一致性支撐。

\subsection*{(三) 會議導向的摘要與組織記憶：從對話摘要到議程／決策導向}
本研究場景（專題會議）本質上屬於多輪、多人、目標驅動的對話序列，與「會議摘要／會議記憶」研究高度相關。早期工作已探討如何從多說話者對話中萃取可讀摘要，並處理冗餘、話輪交錯與資訊重要性判定等問題 \citep{zechner2001concise_dialogue_summaries,wang-cardie-2011-summarizing}。近期工作進一步將\emph{議程（agenda）}、\emph{決策點}與\emph{行動項（action items）}視為會議摘要的關鍵結構，強化「可追蹤的會議脈絡」而非僅生成流暢文字 \citep{liu2025agendaaware_meetingsum,laskar-etal-2023-building}。這些研究共同指出：會議資訊的價值在於其可執行性與可追溯性，特別是「決策理由」與「後續行動」對長程任務影響最大。

對「專題指導一致性」而言，會議摘要若只生成自然語言段落，仍難以穩定支援跨會議推理：同一 TODO 的前置條件、方法修改的理由、實驗失敗的原因與替代方案等，往往需要以結構化關係被保存，才能避免後續指導在不同會議間互相矛盾。故本研究以「Dependency Graph」作為長期記憶載體，將會議輸入由摘要提升為\emph{可運算、可檢核}的研究依賴表示。

\subsection*{(四) 圖式／結構化長期記憶：以 Dependency Graph 支撐跨會議推理}
圖式記憶近年在 LLM 代理領域快速發展，其核心概念是讓系統以節點／邊表示「事實、關係、因果、任務分解」並支援精準檢索與一致性檢查。GraphAgentMemory 強調以圖結構組織代理的經驗與知識，可提升多步推理與長期任務的穩健度 \citep{a2602_05665v1_graphagentmemory}。此類方法與傳統「僅以向量檢索記憶片段」相比，關鍵優勢在於：圖能顯式表示依賴鏈與衝突檢查規則，使跨會議的一致性可以被系統性地檢測（例如：若節點 A 依賴 B，但 B 尚未完成，則禁止輸出與 A 完成相矛盾的結論），從而降低跨會議推理中產生邏輯矛盾的可能性。

本研究採「多層次記憶協調」：Working/Short-term Memory 著重即時對話與近程 TODO；Long-term Memory 以 Dependency Graph 保存研究脈絡（目標、假設、方法、實驗、結果、決策、TODO）及其依賴關係。這樣的設計使「一致性」不再僅靠模型自發保持，而是由外部可檢核結構提供約束與提示，形成可落地的系統化可靠性路徑 \citep{a2310_08560v2_memgpt,a2602_05665v1_graphagentmemory}。

\subsection*{(五) Persona Drift 與可控人設：從角色設計到表徵監控}
第二個長跨度失效在於 persona drift：在長期互動中，模型可能受使用者回覆影響而偏移為迎合式、逢迎式回應，或失去學術嚴謹的指導風格。教育對話研究早已指出「角色」與「回應意圖」會顯著影響學習成效與互動品質；例如對話意圖與角色策略的建模，被用於使系統在不同教學階段產生適切回應 \citep{petukhova-kochmar-2025-intent,tao-etal-2024-rolecraft}。然而，傳統角色設計多透過 prompt 或資料蒐集維持，較難在長期對話中提供\emph{可量測、可干預}的穩定性保證。

Persona Vectors 的研究提出：人設特質可在模型內部表徵空間中以方向向量近似，並可用於監測與引導生成行為，提供「\emph{以表徵為介面}」的控制途徑 \citep{a2507_21509v3_personavectors}。此一觀點對本研究至關重要：我們不僅要「定義教授人設」，更要能量化當前生成是否仍對齊教授式指導特質，並在偏移時採取即時校正。

\subsection*{(六) 推理時干預與表徵工程：以 Residual Stream 操控提升可靠性}
在不進行高成本 retraining 的前提下，「推理時干預」提供一條輕量、可部署的可靠性策略。Activation Engineering 類方法展示可透過在激活空間加入特定方向來 steer 模型行為，並以較低代價達成顯著控制效果 \citep{turner-etal-2023-activation-addition}。Representation Engineering（RepE）綜述更系統地整理「辨識表徵—操作化—控制」的流程與挑戰，指出表徵操控在多概念管理、可靠性與效能保持方面仍需嚴謹設計 \citep{wehner-etal-2025-representation-engineering}。

在更廣義的「代理式（agentic）」系統設計上，ReAct 以「推理—行動」交替的格式化軌跡，證明語言模型可在思考過程中適時呼叫工具或外部知識，以提升任務完成率與可解釋性 \citep{Yao2023ReAct}；Reflexion 進一步引入自我反思與經驗記錄，讓代理能在失敗後以自然語言回饋形成更穩定的行為策略 \citep{Shinn2023Reflexion}。此外，Generative Agents 顯示多代理的日常行為可由記憶、反思與規劃（planning）驅動而產生可持續互動 \citep{Park2023Generative}。上述研究為本計畫的系統化設計提供補強：除記憶與表徵操控外，我們亦可在虛擬教授的推理流程中納入「反思—修正」迴路與可追溯的行動記錄，使跨會議指導更具可檢核性與可重現性。
綜合上述表徵工程與代理式系統的進展，本計畫據此設計「即時人設向量監控（以內積投影量測）」與「自適應 steering（動態調 $\alpha$）」：以內積分數作為一致性量測，再於中後層殘差流注入教授人設向量進行校正，使 persona drift 可被監測並被即時抑制。

\subsection*{(七) 小結：研究定位與貢獻缺口}
綜合上述，現有文獻在各面向均已累積重要基礎：教育導師評測框架與 benchmark \citep{maurya-etal-2025-unifying,a2510_02663v1_tutorbench,a2601_21375v1_teachbench}、長期記憶系統與管理策略 \citep{a2305_10250v3_memorybank,a2310_08560v2_memgpt,a2504_19413v1_mem0}、會議摘要與決策導向表示 \citep{liu2025agendaaware_meetingsum,laskar-etal-2023-building}，以及 persona 表徵與推理時控制 \citep{a2507_21509v3_personavectors,turner-etal-2023-activation-addition,wehner-etal-2025-representation-engineering}。
但這些脈絡尚未被整合，尚缺少一個針對「跨會議學術指導」的整合式系統：同時以圖式長期記憶確保研究脈絡一致，並以表徵層的人設監控＋推理時干預來抑制 persona drift。本研究即以此整合缺口為目標，提出可實作的架構與可量測的評估設計，作為高可靠性教育代理之方法論參考。